\documentclass[12pt,a4paper]{article}
\usepackage[UTF8]{ctex}
\usepackage{geometry}
\usepackage{listings}
\usepackage{xcolor}
\usepackage{hyperref}
\usepackage{graphicx}

\geometry{margin=2.5cm}

\lstset{
    language=Python,
    basicstyle=\ttfamily\small,
    keywordstyle=\color{blue}\bfseries,
    commentstyle=\color{green!60!black},
    stringstyle=\color{red},
    numbers=left,
    numberstyle=\tiny\color{gray},
    stepnumber=1,
    numbersep=5pt,
    frame=single,
    breaklines=true,
    showstringspaces=false,
    tabsize=4
}

\title{Halo 系统代码文档：完整技术文档}
\author{代码分析文档}
\date{\today}

\begin{document}

\maketitle

\tableofcontents
\newpage

\section{文档概述}

本文档是 Halo 系统的完整技术文档，详细分析了系统的各个模块，从函数级别到语句级别解释了代码的作用和实现逻辑。

\subsection{文档结构}

本文档由以下子文档组成：

\begin{enumerate}
    \item \textbf{Components 模块}：核心数据结构（Operator、Query、ModelConfig、ExecuteInfo）
    \item \textbf{Parser 模块}：YAML 配置解析和 DAG 构建
    \item \textbf{Schedulers 模块}：多种调度策略的实现
    \item \textbf{Optimizers 模块}：多进程协调器和执行引擎
    \item \textbf{Workers 模块}：模型推理执行层
    \item \textbf{Util 模块}：工具函数
\end{enumerate}

\subsection{系统架构}

Halo 系统采用分层架构：

\begin{enumerate}
    \item \textbf{配置层}：YAML 文件定义工作流
    \item \textbf{解析层}：Parser 将配置转换为 DAG
    \item \textbf{调度层}：Scheduler 生成执行计划
    \item \textbf{协调层}：Optimizer 管理多进程执行
    \item \textbf{执行层}：Worker 执行实际推理
\end{enumerate}

\section{系统工作流程}

\subsection{整体流程}

\begin{enumerate}
    \item \textbf{配置加载}
    \begin{itemize}
        \item 从 YAML 文件加载工作流定义
        \item 包含 OP 定义、依赖关系、模型配置等
    \end{itemize}
    
    \item \textbf{图构建}
    \begin{itemize}
        \item Parser 解析配置，构建 Operator 对象图
        \item 验证依赖关系，计算最长距离
        \item 检测循环依赖
    \end{itemize}
    
    \item \textbf{调度规划}
    \begin{itemize}
        \item 根据选定的调度策略，生成执行计划
        \item 将 OP 分配到多个 GPU 设备
        \item 考虑依赖关系、模型切换成本、负载均衡等
    \end{itemize}
    
    \item \textbf{执行协调}
    \begin{itemize}
        \item Optimizer 创建多个 Worker 进程
        \item 按计划派发任务到各 Worker
        \item 检查依赖关系，确保正确执行顺序
        \item 收集结果并更新 Query 状态
    \end{itemize}
    
    \item \textbf{模型推理}
    \begin{itemize}
        \item Worker 接收任务，加载模型（如需要）
        \item 执行批量推理
        \item 返回结果给 Optimizer
    \end{itemize}
    
    \item \textbf{结果收集}
    \begin{itemize}
        \item Optimizer 收集所有 Worker 的结果
        \item 更新 Query 的 op\_output 字段
        \item 计算性能指标
    \end{itemize}
\end{enumerate}

\subsection{数据流}

\begin{verbatim}
YAML 配置
    ↓
Parser → Operator 图
    ↓
Scheduler → Workflows (执行计划)
    ↓
Optimizer → ExecuteInfo (任务封装)
    ↓
Worker → 推理结果
    ↓
Optimizer → 更新 Query.op_output
    ↓
最终结果
\end{verbatim}

\section{关键设计决策}

\subsection{多进程架构}

\begin{itemize}
    \item \textbf{原因}：Python GIL 限制，多进程可以真正并行利用多 GPU
    \item \textbf{实现}：每个 GPU 对应一个独立的 Worker 进程
    \item \textbf{通信}：使用 \texttt{multiprocessing.Queue} 进行进程间通信
\end{itemize}

\subsection{依赖关系保证}

\begin{itemize}
    \item \textbf{问题}：工作流中存在依赖关系，必须保证执行顺序
    \item \textbf{解决}：
    \begin{enumerate}
        \item 调度器按拓扑顺序分配任务
        \item Optimizer 在执行前检查依赖是否满足
        \item 不满足依赖的任务会被阻塞，等待父任务完成
    \end{enumerate}
\end{itemize}

\subsection{批量处理优化}

\begin{itemize}
    \item \textbf{策略}：将多个查询打包成 batch，一次推理处理
    \item \textbf{优势}：提高 GPU 利用率，减少推理延迟
    \item \textbf{实现}：vLLM 自动管理动态批处理
\end{itemize}

\subsection{模型切换优化}

\begin{itemize}
    \item \textbf{问题}：切换模型需要重新加载，成本高
    \item \textbf{解决}：
    \begin{enumerate}
        \item 调度器尽量将相同模型的 OP 分配到同一设备
        \item Beam-Search 调度器考虑模型切换成本
        \item Worker 延迟加载模型，只在需要时切换
    \end{enumerate}
\end{itemize}

\section{各模块详细文档}

\subsection{Components 模块}

\textbf{文件}：\texttt{halo/components/*.py}

\textbf{主要内容}：
\begin{itemize}
    \item \texttt{Operator}：工作流节点，包含模型配置和依赖关系
    \item \texttt{Query}：用户查询，维护执行状态和中间结果
    \item \texttt{ModelConfig}：模型配置参数封装
    \item \texttt{ExecuteInfo}：执行任务信息封装
\end{itemize}

\textbf{详细文档}：参见 \texttt{components.tex}

\subsection{Parser 模块}

\textbf{文件}：\texttt{halo/parser.py}

\textbf{主要内容}：
\begin{itemize}
    \item YAML 配置解析
    \item Operator 图构建
    \item 依赖关系建立
    \item 最长距离计算
    \item 循环依赖检测
\end{itemize}

\textbf{详细文档}：参见 \texttt{parser.tex}

\subsection{Schedulers 模块}

\textbf{文件}：\texttt{halo/schedulers/*.py}

\textbf{主要内容}：
\begin{itemize}
    \item \texttt{schedule\_rr}：轮询调度
    \item \texttt{schedule\_by\_levels}：层级调度
    \item \texttt{schedule\_colocate\_parents}：父节点共置调度
    \item \texttt{schedule\_dp}：数据并行调度
    \item \texttt{schedule\_heuristic}：启发式调度
    \item \texttt{schedule\_search}：Beam-Search 调度
\end{itemize}

\textbf{详细文档}：参见 \texttt{schedulers.tex}

\subsection{Optimizers 模块}

\textbf{文件}：\texttt{halo/optimizers/halo\_v.py}

\textbf{主要内容}：
\begin{itemize}
    \item Worker 进程管理
    \item 调度器调用
    \item 任务派发和结果收集
    \item 依赖关系检查
    \item 性能监控
\end{itemize}

\textbf{详细文档}：参见 \texttt{optimizers.tex}

\subsection{Workers 模块}

\textbf{文件}：\texttt{halo/workers/worker\_v.py}

\textbf{主要内容}：
\begin{itemize}
    \item 模型加载和管理
    \item 批量推理执行
    \item 结果格式化
    \item 命令循环处理
\end{itemize}

\textbf{详细文档}：参见 \texttt{workers.tex}

\subsection{Util 模块}

\textbf{文件}：\texttt{halo/util.py}

\textbf{主要内容}：
\begin{itemize}
    \item 数据类型转换
    \item GPU 设备管理
\end{itemize}

\textbf{详细文档}：参见 \texttt{util.tex}

\section{使用示例}

\subsection{基本使用}

\begin{lstlisting}
from halo.optimizers import Optimizer
from halo.components import Query

# 1. 创建优化器（加载配置）
opt = Optimizer("templates/adv_reason_3.yaml")

# 2. 创建查询
queries = [
    Query(0, "What is Machine Learning?"),
    Query(1, "Explain neural networks."),
]

# 3. 调度和执行
opt.schedule(queries, strategy="heuristic")
queries = opt.execute(return_queries=True)

# 4. 查看结果
for q in queries:
    print(f"Query {q.id}:")
    for op_id, output in q.op_output.items():
        print(f"  {op_id}: {output[:100]}...")

# 5. 性能分析
opt.print_latency_percentiles()

# 6. 清理
opt.exit()
\end{lstlisting}

\subsection{自定义工作流}

创建 YAML 配置文件：

\begin{lstlisting}
start_ops:
  - op0
end_ops:
  - op2

ops:
  op0:
    model: meta-llama/Llama-3.1-8B-Instruct
    prompt: "Answer the question."
    max_tokens: 512
    input_ops: []
    output_ops:
      - op1
  
  op1:
    model: meta-llama/Llama-3.1-8B-Instruct
    prompt: "Refine the answer."
    max_tokens: 1024
    input_ops:
      - op0
    output_ops:
      - op2
  
  op2:
    model: meta-llama/Llama-3.1-8B-Instruct
    prompt: "Provide final answer."
    max_tokens: 512
    input_ops:
      - op1
    output_ops: []
\end{lstlisting}

\section{性能优化建议}

\subsection{调度策略选择}

\begin{itemize}
    \item \textbf{简单工作流}：使用 \texttt{rr} 或 \texttt{by\_levels}
    \item \textbf{多模型切换}：使用 \texttt{search} 或 \texttt{heuristic}
    \item \textbf{大批量查询}：使用 \texttt{dp}（数据并行）
    \item \textbf{缓存优化}：使用 \texttt{colocate\_parents}
\end{itemize}

\subsection{配置优化}

\begin{itemize}
    \item \textbf{批处理大小}：根据 GPU 显存调整 \texttt{max\_batch\_size}
    \item \textbf{模型量化}：使用量化模型减少显存占用
    \item \textbf{keep\_cache}：对于相同模型的连续 OP，启用 cache 保留
\end{itemize}

\section{总结}

Halo 系统是一个完整的智能体工作流执行框架，具有以下特点：

\begin{itemize}
    \item \textbf{声明式配置}：通过 YAML 文件定义工作流，易于修改和扩展
    \item \textbf{多设备并行}：支持多 GPU 并行执行，提高吞吐量
    \item \textbf{智能调度}：多种调度策略，优化执行效率
    \item \textbf{依赖保证}：自动检查依赖关系，确保正确执行
    \item \textbf{高性能}：基于 vLLM，支持高效的批量推理
\end{itemize}

通过模块化的设计和清晰的接口，Halo 系统实现了灵活性和性能的平衡。

\section{附录}

\subsection{相关文档}

\begin{itemize}
    \item \texttt{components.tex}：Components 模块详细文档
    \item \texttt{parser.tex}：Parser 模块详细文档
    \item \texttt{schedulers.tex}：Schedulers 模块详细文档
    \item \texttt{optimizers.tex}：Optimizers 模块详细文档
    \item \texttt{workers.tex}：Workers 模块详细文档
    \item \texttt{util.tex}：Util 模块详细文档
\end{itemize}

\subsection{代码结构}

\begin{verbatim}
halo/
├── components/      # 核心数据结构
├── optimizers/      # 优化器和协调器
├── parser.py        # 配置解析
├── schedulers/      # 调度策略
├── workers/         # 执行层
└── util.py          # 工具函数
\end{verbatim}

\end{document}
